\section{线性 ODE 与精确解}

\subsection{简单线性 ODE 的解析解}

在讨论扩散模型的 ODE 之前，讲者从最简单的线性 ODE 入手建立直觉。考虑标量 ODE：
$$\frac{dx}{dt} = p(t) \cdot x(t)$$

其中 $p(t)$ 是已知的时间函数。假设 $x(t) > 0$，两边除以 $x(t)$：
$$\frac{1}{x(t)} \frac{dx}{dt} = p(t)$$

左侧恰好是 $\frac{d}{dt} \ln x(t)$，因此：
$$\ln x(t) - \ln x(s) = \int_s^t p(\tau)\, d\tau$$

\begin{importantbox}{线性 ODE 的精确解}
$$x(t) = x(s) \cdot \exp\left(\int_s^t p(\tau)\, d\tau\right)$$
这个解是\textbf{解析}的——不需要任何离散化近似。对于线性 ODE，我们总能通过积分得到精确解。这一事实是后续推导的关键基石。
\end{importantbox}

\subsection{半线性 ODE 与积分因子法}

扩散模型的 PF-ODE 不是纯线性的，而是\textbf{半线性}（semi-linear）的，形如：
$$\frac{d\mathbf{x}}{dt} = p(t) \cdot \mathbf{x}(t) + q(t)$$

其中 $p(t) \cdot \mathbf{x}$ 是线性部分，$q(t)$ 是非线性偏移项（在扩散模型中，$q(t)$ 包含噪声预测网络的输出）。

为了求解半线性 ODE，引入\textbf{积分因子}（integrating factor）：
$$\mu(t) = \exp\left(-\int_0^t p(\tau)\, d\tau\right)$$

$\mu(t)$ 是齐次方程 $\frac{dx}{dt} = p(t) \cdot x$ 的"逆解"。将 ODE 两边乘以 $\mu(t)$：
$$\mu(t) \frac{dx}{dt} - \mu(t) p(t) x(t) = \mu(t) q(t)$$

利用乘积求导法则，左侧化为 $\frac{d}{dt}\left[\mu(t) x(t)\right]$，积分后得到：

\begin{importantbox}{半线性 ODE 的精确解（Variation of Constants）}
$$x(t) = \frac{\mu(s)}{\mu(t)} x(s) + \frac{1}{\mu(t)} \int_s^t \mu(\tau) q(\tau)\, d\tau$$
\begin{itemize}
    \item 第一项：线性部分的精确解，只依赖初始值 $x(s)$，可以解析计算
    \item 第二项：非线性部分的积分，涉及 $q(\tau)$（包含神经网络输出），通常需要数值近似
\end{itemize}
这就是\textbf{指数积分器}（Exponential Integrator）的基本形式：线性部分精确处理，仅对非线性部分做数值近似。
\end{importantbox}

\begin{knowledgebox}{为什么叫"指数积分器"？}
因为线性部分的精确解包含指数函数 $\exp(\cdot)$。传统欧拉法会对整个 ODE（线性+非线性）做近似，而指数积分器只对非线性部分做近似。在线性部分占主导的情况下（如扩散模型的 PF-ODE），这种拆分能显著提高精度。
\end{knowledgebox}

\subsection{带常数偏移项的精确解}

作为练习，考虑带常数 $q$（不依赖 $t$）的情形：
$$\frac{dx}{dt} = p(t) \cdot x(t) + c$$

将通解公式代入并化简，可以得到一个更紧凑的表达式。讲者在课上带学生推导了这一过程，核心思路是将 $c$ 提出积分号外，然后利用 $\mu(t)$ 的性质化简。

\subsection{本章小结}

线性 ODE 有解析解；半线性 ODE 可以通过积分因子法将线性部分精确求解，仅剩非线性部分需要数值近似。这是 DPM-Solver 的数学基础。


